\documentclass[11pt]{article}

\usepackage[margin=1in]{geometry}
\usepackage{enumitem}
\usepackage{amsmath, amssymb}
\usepackage{hyperref}

\setlist[enumerate]{itemsep=0.8em, topsep=0.8em}
\setlist[itemize]{itemsep=0.3em, topsep=0.3em}

\begin{document}

\begin{center}
  {\Large \textbf{MATH 340 Examination}}\\[0.25em]
  {\large \textbf{Multivariable Calculus and Vector Analysis}}\\[0.5em]
\end{center}

\noindent\textbf{Instructions:}
\begin{itemize}
  \item Answer all questions.
  \item For Questions 1--4, choose the best option.
  \item For Questions 5--8, write brief but precise answers.
\end{itemize}

\vspace{0.5em}

\begin{enumerate}[label=\textbf{\arabic*.}]

  \item What geometric property does the gradient vector $\nabla f$ at a point represent for a 
  differentiable function $f: \mathbb{R}^n \to \mathbb{R}$?
  \begin{enumerate}[label=(\Alph*)]
    \item The direction of maximum decrease of $f$
    \item A vector tangent to the level set of $f$ at that point
    \item The direction of steepest ascent and perpendicular to level sets
    \item The curvature of the surface defined by $f$
  \end{enumerate}

  \item Which of the following conditions is sufficient for a function $f: \mathbb{R}^2 \to \mathbb{R}$ 
  to have its mixed partial derivatives equal, i.e., $f_{xy} = f_{yx}$?
  \begin{enumerate}[label=(\Alph*)]
    \item $f$ is differentiable at the point
    \item $f$ has continuous first partial derivatives
    \item $f$ has continuous second partial derivatives
    \item $f$ is twice differentiable at the point
  \end{enumerate}

  \item The Jacobian determinant of a transformation $T: \mathbb{R}^n \to \mathbb{R}^n$ at a point 
  provides information about:
  \begin{enumerate}[label=(\Alph*)]
    \item The rate of change of $T$ along coordinate axes
    \item The local scaling factor of volumes under the transformation
    \item Whether $T$ is injective globally
    \item The total derivative matrix of $T$
  \end{enumerate}

  \item Which statement correctly describes the relationship between Green's Theorem and Stokes' Theorem?
  \begin{enumerate}[label=(\Alph*)]
    \item Green's Theorem is a special case of Stokes' Theorem in two dimensions
    \item Stokes' Theorem generalizes the Divergence Theorem to higher dimensions
    \item Green's Theorem applies to simply connected regions only
    \item They are independent theorems with no direct relationship
  \end{enumerate}

  \item Explain the geometric and computational significance of the Hessian matrix in analyzing 
  critical points of a function $f: \mathbb{R}^n \to \mathbb{R}$. How does the second derivative 
  test for functions of several variables use the Hessian?

  \item Describe the difference between a conservative vector field and a path-independent vector 
  field. Under what conditions are these properties equivalent, and what role does the curl play?

  \item State the Divergence Theorem (Gauss's Theorem) and explain its physical interpretation in 
  the context of fluid flow. How does it relate local behavior (divergence) to global behavior (flux)?

  \item Explain why changing the order of integration in a double integral requires careful 
  consideration of the region of integration. Provide a method for determining the new limits 
  when switching from $\int \int dx\,dy$ to $\int \int dy\,dx$.

\end{enumerate}

\end{document}
